%-------------------------------------------------------------------------------
%	SECTION TITLE
%-------------------------------------------------------------------------------
\cvsection{Publications}
\begin{cvpubs}

    \cvpub{2026}


    \cvpub{[19] Dunnigan, A. H., \textbf{Liu, D.}, Steinbrügge, G. B., Rivoldini, A., Dumberry, M., Cao, H., Soderlund, K. M.: Interior models of Mercury and conditions for iron snow formation in a Fe-S-Si core. \textit{J. Geophys. Res. - Planets}, 131, e2025JE009368, 2026.}


    \cvpub{2025}

    \cvpub{[18] \textbf{Liu, D.}, Becker, TW.: Earthquake Rupture Dynamics From Graph Neural Networks. \textit{J. Geophys. Res. - Solid Earth}, 130(12), e2025JB031981, 2025.}

    \cvpub{[17] White, K., \textbf{Liu, D.}, Persad, G.: Absence of Aerosol Indirect Effect Dependence on Background Climate State in NCAR CESM2. \textit{J. Climate}, \textbf{38}, 147-163, https://doi.org/10.1175/JCLI-D-23-0755.1, 2025.}

    \cvpub{2024}

    \cvpub{[16] Thirumalai, K., DiNezio, PN., Partin, JW., \textbf{Liu, D.}, Costa, K., and Jacobel, A.: Future increase in extreme El Niño supported by past glacial changes. \textit{Nature}, https://doi.org/10.1038/s41586-024-07984-y, 2024.}

    \cvpub{[15] \textbf{Liu, D.}, Puel, S., Becker, TW., Moresi, L.: Analytical and numerical models of viscous anisotropy: a toolset to constrain the role of mechanical anisotropy for regional tectonics and fault loading. \textit{Geophys. J. Int.}, 239(2), 950-963, https://doi.org/10.1093/gji/ggae296, 2024.}

    \cvpub{[14] Puel, S., Becker, TW., Villa, U., Ghattas, O., \textbf{Liu, D.}: Volcanic arc rigidity variations illuminated by coseismic deformation of the 2011 Tohoku-oki M9. \textit{Science Adv.}, 10, eadl4264, https://doi.org/10.1126/sciadv.adl4264, 2024.}

    \cvpub{[13] Puel, S., Becker, TW., Villa, U., Ghattas, O., \textbf{Liu, D.}: An adjoint-based optimization method for jointly inverting heterogeneous material properties and fault slip from earthquake surface deformation data. \textit{Geophys. J. Int.}, 236(2), 778-797, https://doi.org/10.1093/gji/ggad442, 2024.}

    \cvpub{2023}

    \cvpub{[12] Xu, X., \textbf{Liu, D.}, Lavier, L.: Constraining fault damage zone properties from geodesy: A case study near the 2019 Ridgecrest earthquake sequence. \textit{Geophys. Res. Lett.}, 50, e2022GL101692, 2023.}

    \cvpub{2022}

    \cvpub{[11] Puel, S., Khattatov, E., Villa, U., \textbf{Liu, D.}, Ghattas, O., Becker TW: A mixed, unified forward/inverse framework for earthquake problems: fault implementation and coseismic slip estimate. \textit{Geophys. J. Int.}, 230(2), 733-758, 2022.}

    \cvpub{[10] Jiang, J. et al.: Community-driven code comparisons for three-dimensional dynamic modeling of sequences of earthquakes and aseismic slip. \textit{J. Geophys. Res. - Solid Earth}, 127(3), e2021JB023519, 2022.}

    \cvpub{[9] \textbf{Liu, D.}, Duan, B., Scharer, K., Yule, D.: Observation-constrained multicycle dynamic models of the southern San Andreas and the northern San Jacinto faults: addressing complexity in paleoearthquake extent and recurrence with realistic 2D fault geometry. \textit{J. Geophys. Res. - Solid Earth}, 127(2), e2021JB023420, 2022.}

    \cvpub{2021}

    \cvpub{[8] \textbf{Liu, D.}, Duan, B., Prush, VB, Oskin, M., Liu-Zeng, J.: Observation-constrained multicycle dynamic models of the Pingding Shan earthquake gate along the Altyn Tagh Fault. \textit{Tectonophys.}, 814, 228948, 2021.}

    \cvpub{2020}

    \cvpub{[7] \textbf{Liu, D.}, Duan, B., Luo, B.: EQsimu, a 3-D finite element dynamic earthquake simulator for multicycle dynamics of geometrically complex faults governed by rate- and state-dependent friction. \textit{Geophys. J. Int.}, 220(1), 598-609, 2020.}

    \cvpub{[6] Luo, B., Duan, B, \textbf{Liu, D.}: 3D Finite element modeling of dynamic rupture and aseismic slip over earthquake cycles on geometrically complex faults. \textit{Bulletin of Seismological Society of America}, 110(6), 2619-2637, 2020.}

    \cvpub{2015-2019}

    \cvpub{[5] Zhong, S., Wan, Z., Duan, B., \textbf{Liu, D.}, Luo, B.: Do earthquakes trigger mud volcanoes? A case study from the Southern margin of the Junggar Basin, NW China. \textit{Geological J.}, 54(3), 1223-1237, 2019.}

    \cvpub{[4] \textbf{Liu, D.}, Duan, B.: Scenario earthquake and ground motion simulations in North China: effects of heterogeneous fault stress and 3D basin structure. \textit{Bulletin of Seismological Society of America}, 108(4), 2148-2169, 2018.}

    \cvpub{[3] Harris, RA, Barall, M., Aagaard, B., Ma, S., Roten, D., Olsen, K., Duan, B, \textbf{Liu, D.}, et al.: A suite of exercises for verifying dynamic earthquake rupture codes. \textit{Seis. Res. Lett.}, 89(3), 1146-1162, 2018.}

    \cvpub{[2] Duan, B., \textbf{Liu, D.}, Yin, A: Seismic shaking in the North China Basin expected from ruptures of a possible seismic gap. \textit{Geophys. Res. Lett.}, 44(10), 4855-4862, https://doi.org/10.1002/2017GL072638, 2017.}

    \cvpub{[1] \textbf{Liu, D.}, Hu, C., Cai, Y.: Numerical simulation of the dynamic rupture process of the 2011 Tohoku-Oki, Japan, Mw9.0 earthquake. \textit{Chinese J. of Geophys.}, 58(9), 3133-3143, 2015 (in Chinese).}
    
\end{cvpubs}

%-------------------------------------------------------------------------------
\cvsubsection{Preprints}
\begin{cvpubs}
    \cvpub{[P4] \textbf{Liu, D.}, You, K., Ren, P., Tang, H.: Nonlinearly coupled permafrost dynamics from graph neural networks. \textit{Research Square}, https://doi.org/10.21203/rs.3.rs-10910110/v1, 2026.}

    \cvpub{[P3] Song, C., Puel, S., \textbf{Liu, D.}, Wallace, L., Becker, TW.: The role of elastic heterogeneity for inferring seismic-cycle deformation at the Nicoya Peninsula, Costa Rica. \textit{ESS Open Archive}, https://doi.org/10.22541/essoar.15005894/v1, 2026.}

    \cvpub{[P2] Yun, J., Fialko, Y., May, D., Gabriel, A.-A., Williams, C., \textbf{Liu, D.}: Effects of loading schemes in volumetric simulations of sequences of earthquakes and aseismic slip (SEAS) in subduction zones. \textit{EarthArXiv}, https://doi.org/10.31223/X5D20T, 2026.}

    \cvpub{[P1] Domeyko, R., Partin, J., Okumura, Y., Thirumalai, K., DiNezio, P., et al. (incl. \textbf{Liu, D.}): Rare deglacial-aged corals point to weaker El Ni\~no in a colder world. \textit{EarthArXiv}, https://doi.org/10.31223/X5R50Z, 2026.}
\end{cvpubs}
% %---------------------------------------------------------
%\cvsubsection{In Prep}
% %---------------------------------------------------------

%\begin{cvpubs}
%\small \color{black}
%  \cvpub{Manuscript 1}

 %   \cvpub{Manuscript 2}
%\end{cvpubs}

% %---------------------------------------------------------
